\chapter{Dengue Fever}

\section*{Definition and Overview}
Dengue fever is the most common form of illness caused by the dengue virus, a mosquito-borne infection that produces a sudden high fever with severe headache, muscle and joint pain, and often a rash. It is the classic presentation of dengue infection and, for most people, is an unpleasant but self-limiting illness lasting about a week. The importance of dengue fever lies in its potential to progress to severe dengue, with bleeding and shock, which is why careful monitoring is essential. This chapter explains the features, course, and care of dengue fever itself.

\section*{Epidemiology}
Dengue fever is widespread across tropical and subtropical regions, with the greatest burden in Southeast Asia, the Western Pacific, South Asia, and the Americas. The World Health Organization estimates hundreds of millions of infections occur annually, with symptomatic dengue fever representing a substantial proportion. In Australia, dengue fever is seen in travellers returning from endemic areas and in localised outbreaks in Far North Queensland. The number of cases fluctuates with rainfall, temperature, and mosquito populations, and has been rising globally as the Aedes mosquito expands its range.

\section*{Aetiology and Pathophysiology}
\begin{itemize}
    \item \textbf{Dengue virus:} one of four serotypes causes the infection; infection is transmitted by the bite of infected Aedes aegypti mosquitoes
    \item \textbf{Viral entry:} the virus enters cells of the immune system and multiplies, spreading through the bloodstream during the febrile phase
    \item \textbf{Immune response:} the immune reaction produces the fever, headache, and muscle pain
    \item \textbf{Serotype immunity:} recovery gives lifelong immunity to the infecting serotype but only partial, temporary protection against others
    \item \textbf{Risk of severe disease:} most first infections cause dengue fever, but subsequent infection with a different serotype carries a higher risk of severe dengue
\end{itemize}

\section*{Clinical Features}
\begin{itemize}
    \item Abrupt onset of high fever, typically 39--40 degrees Celsius
    \item Severe headache, especially behind the eyes (retro-orbital pain)
    \item Severe muscle and joint aches ("break-bone fever")
    \item Nausea, vomiting, and poor appetite
    \item A flushed or blotchy rash, which may appear a few days into the illness
    \item Swollen glands in some people
    \item The fever typically lasts 2--7 days
    \item Warning signs that may herald severe disease: severe abdominal pain, persistent vomiting, bleeding gums or nose, drowsiness, and restlessness
\end{itemize}

\section*{Differential Diagnosis}
\begin{itemize}
    \item Other mosquito-borne infections, including Zika, chikungunya, and Ross River virus
    \item Malaria, in travellers from affected regions
    \item Influenza, COVID-19, and other viral respiratory infections
    \item Typhoid fever and other enteric fevers
    \item Measles and rubella, which cause fever with rash
    \item Acute gastroenteritis, when vomiting is prominent
\end{itemize}

\section*{When to Seek Medical Care}
Anyone who develops a fever within two weeks of travelling to a dengue-affected area should see a doctor and report their travel history. People with suspected or confirmed dengue need monitoring, and medical review is required if symptoms worsen or warning signs appear, especially as the fever settles.

\textbf{Call 000 (or local emergency services) for:}
\begin{itemize}
    \item Severe abdominal pain, persistent vomiting, or bleeding from gums, nose, or bowel
    \item Drowsiness, confusion, or fainting
    \item Difficulty breathing, cold skin, or a weak rapid pulse (possible shock)
    \item Inability to drink or keep down fluids
\end{itemize}

\section*{Diagnosis and Investigations}
\begin{itemize}
    \item \textbf{Clinical diagnosis:} fever with headache, muscle pain, and rash in a person with relevant exposure
    \item \textbf{PCR testing:} detects the virus in blood during the first week of illness
    \item \textbf{NS1 antigen test:} a rapid test positive early in the illness
    \item \textbf{Serology:} IgM and IgG antibodies appear from around day 4--5 and confirm recent infection
    \item \textbf{Full blood count:} monitoring platelets and haematocrit to detect progression to severe dengue
    \item \textbf{Liver and other blood tests:} as clinically indicated
\end{itemize}

\section*{Management}
\begin{itemize}
    \item \textbf{Home care for mild cases:} rest, fluids, and paracetamol for fever and pain
    \item \textbf{Medication caution:} avoid aspirin, ibuprofen, and other non-steroidal anti-inflammatory drugs, which increase bleeding risk
    \item \textbf{Hydration:} oral rehydration solutions and plenty of fluids; intravenous fluids if the person cannot drink
    \item \textbf{Monitoring:} daily review and blood counts during the critical period
    \item \textbf{Hospital admission:} for warning signs, co-existing conditions, pregnancy, or inability to manage at home
    \item \textbf{Treatment of severe features:} hospital care with careful fluid therapy and blood products if needed
    \item \textbf{Follow-up:} review after recovery for persisting fatigue or complications
\end{itemize}

\section*{Complications}
\begin{itemize}
    \item Progression to severe dengue with plasma leakage, bleeding, and shock
    \item Dehydration from fever, vomiting, and poor intake
    \item Prolonged fatigue and muscle weakness after the acute illness
    \item Rare complications such as liver involvement or encephalopathy
    \item Complications in pregnancy, including increased risk of preterm birth
\end{itemize}

\section*{Prognosis and Outcomes}
The vast majority of people with dengue fever recover fully within one to two weeks with rest and fluids, and the illness is self-limiting. The key to a good outcome is recognising the small proportion of people who progress to severe dengue, particularly around the time the fever breaks, and getting them to hospital promptly. With good supportive care, severe dengue is treatable, and death is rare. Immunity to the infecting serotype is lifelong, but travellers and residents in endemic areas remain at risk from other serotypes.

\section*{Prevention and Self-Care}
\begin{itemize}
    \item Use effective insect repellent and wear protective clothing in dengue-affected areas
    \item Remove standing water where mosquitoes breed, such as in pots, tyres, and containers
    \item Take bite precautions both day and night, as Aedes mosquitoes bite during the day
    \item Report fever after travel to a doctor promptly
    \item Rest, drink fluids, and use paracetamol rather than aspirin during illness
    \item Seek urgent review if warning signs of severe disease appear
\end{itemize}

\section*{Sources and Further Reading}
The following reputable sources provide further information for healthcare students and patients seeking reliable, current advice about dengue fever.

\begin{itemize}
    \item World Health Organization. \textit{Dengue and severe dengue.} https://www.who.int/
    \item Centers for Disease Control and Prevention. \textit{Dengue fever: clinical information.} https://www.cdc.gov/dengue/
    \item Queensland Health. \textit{Dengue fever: symptoms and treatment.} https://www.health.qld.gov.au/
    \item NHS. \textit{Dengue: symptoms and treatment.} https://www.nhs.uk/conditions/dengue/
    \item MedlinePlus. \textit{Dengue fever.} https://medlineplus.gov/dengue.html
    \item Healthdirect Australia. \textit{Dengue fever.} https://www.healthdirect.gov.au/
\end{itemize}
